\documentclass[a4paper,11pt]{article}
\usepackage[utf8]{inputenc}
\usepackage[T1]{fontenc}
\usepackage[french]{babel}
\usepackage[left=2.5cm,right=2.5cm,top=1.8cm,bottom=1.8cm]{geometry}
\usepackage{hyperref}
\usepackage{xcolor}
\usepackage{fontawesome5}
\usepackage{lmodern}

\definecolor{primary}{RGB}{0, 82, 147}
\hypersetup{colorlinks=true, linkcolor=primary, urlcolor=primary}

\begin{document}

\noindent
\textbf{Loïc Aron MBASSI EWOLO} \\
Élève-Ingénieur ENSEA / Polytechnique Yaoundé \\
Cergy, France \\
\faPhone\ (+33) 7 59 52 33 98 \quad \faEnvelope\ \href{mailto:loicaron.mbassiewolo@ensea.fr}{loicaron.mbassiewolo@ensea.fr}

\vspace{0.6cm}

\noindent
\textbf{CEA DAM Île-de-France} \\
À l'attention de M. Arnaud PARENT et M. Paul CHAPELLIER \\
Bruyères-le-Châtel - 91297 Arpajon

\hfill Cergy, le 23 janvier 2026

\vspace{0.5cm}

\noindent\textbf{Objet : Stage - Caractérisation de capteurs inertiels MEMS}

\vspace{0.4cm}

Messieurs,

\vspace{0.3cm}

Je suis en deuxième année à l'ENSEA, en double diplôme avec Polytechnique Yaoundé. Je cherche un stage de 4 à 6 mois, et votre offre sur la caractérisation de capteurs MEMS m'intéresse particulièrement.

\vspace{0.25cm}

J'ai déjà travaillé avec des capteurs inertiels. Pour un projet de traduction de langue des signes, j'ai participé à la conception d'un gant instrumenté équipé de capteurs IMU (accéléromètres, gyroscopes) et de flex sensors. Mon rôle incluait la calibration des capteurs, l'acquisition des données et l'analyse de leur comportement. J'ai aussi soudé les composants et intégré le tout sur un microcontrôleur. Ce n'était pas des MEMS de quelques micromètres, mais le processus de caractérisation (comprendre le capteur, le tester, analyser les résultats) est comparable.

\vspace{0.25cm}

À l'ENSEA, je suis la filière électronique et systèmes embarqués. J'ai des cours d'électronique analogique et numérique, d'instrumentation et de mécanique. Pour la programmation, je code principalement en Python. Je n'ai jamais utilisé LabVIEW, mais je suis prêt à l'apprendre si nécessaire.

\vspace{0.25cm}

La mission mentionne l'amélioration du banc de mesure (mécanique, électronique, logiciel). C'est le type de travail que j'apprécie : partir d'un existant, identifier ce qui peut être amélioré, et proposer des modifications. Au laboratoire IoT de mon école, j'ai eu l'occasion de prendre en main des bancs de test et de les adapter à mes besoins expérimentaux.

\vspace{0.25cm}

Concernant la rédaction du rapport de synthèse, j'ai l'habitude de documenter mes travaux. Au Mila (Québec), j'ai rédigé des rapports d'expériences avec analyse statistique des résultats.

\vspace{0.25cm}

Je suis basé à Cergy, Bruyères-le-Châtel est donc accessible. Je suis disponible à partir d'avril 2026.

\vspace{0.4cm}

Cordialement,

\vspace{0.8cm}

\textbf{Loïc Aron MBASSI EWOLO}

\end{document}
